\scanpage{46}
% archive sheet number: 45
\begingroup\csname tagsleft@true\endcsname
\(\cdot\) La donnée de \(\Delta_1\) est unique à conjugaison locale près -- si on prend une trivialisation de \deleted{\unsure{droite de \(X\) sur}} \(\Delta_1\), \deleted{on peut prendre} la situation \((X,\Delta_1)\) i.e. \deleted{\unsure{se}} est isom. localement à \((\mathbb P^1_S,\{0,\infty\})\), \(\sigma_1\) étant donné par \(z\mapsto-z\), donc \(X\setminus\Delta_1\) étant \(\deleted{\unsure{G_m}}\ \mathbb P^1_S\setminus\{0,\infty\}=\mathbb G_{mS}\), et le quotient \((X\setminus\Delta_1)/\{1,\sigma_1\}\), s'identifiant à \(\mathbb G_{mS}\), l'hom de passage au quotient étant
\diagram{\begin{tikzcd}[ampersand replacement=\&,row sep=.15em]
X\setminus\Delta_1\arrow[r]\&(X\setminus\Delta_1)/\{1,\sigma_1\}\\
\rotatebox{90}{\(\simeq\)}\&\rotatebox{90}{\(\simeq\)}\\
\mathbb G_{mS}\arrow[r,"\lambda\mapsto\lambda^2"]\&\mathbb G_{mS}
\end{tikzcd}}

La donnée d'une section de \((X\setminus\Delta_1)/\{1,\sigma_1\}\) équivaut alors à celle d'une section \(\lambda\) de \(\mathbb G_{mS}\) -- et loc pour la top. étale, on peut poser \(\lambda=\unsure{\zeta|\xi}^2\), \deleted{et quitte à remplacer} \inserted{i.e. \deleted{\unsure{que}} la section} \deleted{\unsure{la}} \(\lambda\) provient d'une section \(\zeta\) de \(X\setminus\Delta_1\) sur \(S\). On peut (quitte à remplacer l'iso. \(X\simeq\mathbb P^1_S\) par un autre, savoir son multiple par \(1/\zeta\)) supposer que \unsure{c'est|\(\zeta\) est|\(\xi\) est} la section \(1\), de sorte que la situation \unsure{standart|standard} est
\[
\left\{\begin{aligned}
\Delta_1&=\{0,\infty\},&\Delta_2&=\{1,-1\}\\
\sigma_1(z)&=-z&\sigma_2(z)&=\frac1z
\end{aligned}\right. \tag*{(4)}
\]
(NB \(\sigma_2\) est l'unique homographie qui échange \(0\) et \(\infty\), et fixe \(1\) et \(-1\)). On trouve alors
\endgroup
